\allowdisplaybreaks
\begin{center}
	\hrule
	\vspace{.4cm}
 \begin{tabular*}{\textwidth}{@{}l@{}|@{\extracolsep{0.6in}}r@{}}%
\parbox{4.25in}{\raggedright{\includegraphics[width=.9\linewidth]{ictp-pwf.pdf}}} &
\parbox[c][]{4in}{{\Large\textbf{Md Akiful Islam Zawad} \par}
                    { Brac University \par}
                    { KHA 224, Progati Sarani, Merul Badda, \par}
                    { Dhaka 1212, Bangladesh \par}
                    { \href{ext.akiful.islam@bracu.ac.bd}{ext.akiful.islam@bracu.ac.bd}} \par}
\end{tabular*}\vspace{.3in}
\definecolor{ceruleanblue}{rgb}{0.16, 0.32, 0.6}
	\LARGE\scshape\textbf{\textcolor{ceruleanblue}{Physics for Bangladesh QFT School}}
\end{center}
\hrule\vspace{.25in}
{\large\textbf{Problem Sheet:}\ \textsc{3} \hspace{\hfill} \large\textbf{Due Date:} \today\\
	\hrule}
 %----------------------------
\paragraph*{Problem Set 3} %\hfill \newline
\\
In the lecture, we saw that the $A_\mu$ field has two degrees of freedom. In the Coulomb gauge, this is easy to see since we have the constraints
\begin{align}
    A_0 = 0 \quad \& \quad \nabla \cdot \mathbf{A} = 0.
\end{align}
The second constraint, in momentum space, gives the condition
\begin{align}
    \mathbf{p} \cdot \boldsymbol{\epsilon}(\mathbf{p}) = 0,\label{eq:question-4}
\end{align}
where $\boldsymbol{\epsilon}(\mathbf{p})$ is called the polarization vector. This means that only the components transverse to the momentum $\mathbf{p}$ will survive. Therefore, we can pick two orthonormal basis vectors $\boldsymbol{\epsilon}_a(\mathbf{p})$ ($a = 1, 2$) satisfying Eq. (\ref{eq:question-4}) and
\begin{align}
    \boldsymbol{\epsilon}_a(\mathbf{p}) \cdot \boldsymbol{\epsilon}_b(\mathbf{p}) = \delta_{ab}.
\end{align}

Let the momentum of the photon be $\mathbf{p} = |\mathbf{p}|
\begin{pmatrix}
    \sin \theta \cos \phi \\
    \sin \theta \sin \phi \\
    \cos \theta
\end{pmatrix}$. This momentum can be obtained from the momentum $\mathbf{p}^\prime =
\begin{pmatrix}
    0 \\
    0 \\
    |\mathbf{p}|
\end{pmatrix}$ by a rotation
\begin{align}
    R(\theta, \phi) =
    \begin{pmatrix}
        \cos \phi & -\sin \phi & 0 \\
        \sin \phi & \cos \phi & 0 \\
        0 & 0 & 1
    \end{pmatrix}
    \begin{pmatrix}
        \cos \theta & 0 & \sin \theta \\
        0 & 1 & 0 \\
        -\sin \theta & 0 & \cos \theta
    \end{pmatrix}.
\end{align}

In the primed frame, one can choose the two polarization basis vectors as
\begin{align}
    \boldsymbol{\epsilon}^\prime_1(\mathbf{p}^\prime) =
    \begin{pmatrix}
        1 \\
        0 \\
        0
    \end{pmatrix} \quad \& \quad \boldsymbol{\epsilon}^\prime_2(\mathbf{p}^\prime) =
    \begin{pmatrix}
        0 \\
        1 \\
        0
    \end{pmatrix}.
\end{align}

\begin{itemize}
    \item[(a)] Obtain the two polarization vectors $\boldsymbol{\epsilon}_1(\mathbf{p})$ and $\boldsymbol{\epsilon}_2(\mathbf{p})$ in the unprimed frame by applying the rotation.

    Let us define $\boldsymbol{\epsilon}_3(\mathbf{p}) = \hat{p}$. One can easily see that $\boldsymbol{\epsilon}_r(\mathbf{p})$ ($r = 1, 2, 3$) forms a complete set of orthonormal basis for three-dimensional momentum space. Therefore, they should obey the completeness relation:
    \begin{align}
        \sum_{r=1}^3 \epsilon_i^r(\mathbf{p}) \epsilon_j^r(\mathbf{p}) = \delta_{ij}.
    \end{align}

    \item[(b)] From this, obtain the completeness relation for the two transverse polarization vectors
    \begin{align}
        \sum_{r=1}^2 \epsilon_i^r(\mathbf{p}) \epsilon_j^r(\mathbf{p}) = \delta_{ij} - \frac{p_i p_j}{|\mathbf{p}|^2}.
    \end{align}
\end{itemize}
\bigskip\bigskip\hline\hline\bigskip
\subsection*{Solution:}
\begin{enumerate}
\item [(a)] We start by finding the two polarization vectors in the unprimed frame by applying the rotation matrix.

The momentum of the photon is given as:
\begin{align*}
    \mathbf{p} = |\mathbf{p}|
    \begin{pmatrix}
        \sin \theta \cos \phi \\
        \sin \theta \sin \phi \\
        \cos \theta
    \end{pmatrix}.
\end{align*}
This momentum can be obtained from the momentum in the primed frame:
\begin{align*}
    \mathbf{p}' =
    \begin{pmatrix}
        0 \\
        0 \\
        |\mathbf{p}|
    \end{pmatrix},
\end{align*}
by a rotation matrix \( R(\theta, \phi) \) which is defined as:
\begin{align*}
    R(\theta, \phi) = 
    \begin{pmatrix}
        \cos \phi & -\sin \phi & 0 \\
        \sin \phi & \cos \phi & 0 \\
        0 & 0 & 1
    \end{pmatrix}
    \begin{pmatrix}
        \cos \theta & 0 & \sin \theta \\
        0 & 1 & 0 \\
        -\sin \theta & 0 & \cos \theta
    \end{pmatrix}.
\end{align*}
In the primed frame, we can choose the polarization basis vectors as:
\begin{align}
    \boldsymbol{\epsilon}'_1(\mathbf{p}') =
    \begin{pmatrix}
        1 \\
        0 \\
        0
    \end{pmatrix}, \quad \boldsymbol{\epsilon}'_2(\mathbf{p}') =
    \begin{pmatrix}
        0 \\
        1 \\
        0
    \end{pmatrix}.
\end{align}
To obtain the polarization vectors in the unprimed frame, we apply the rotation matrix \( R(\theta, \phi) \) to these vectors:
\begin{align}
    \boldsymbol{\epsilon}_1(\mathbf{p}) &= R(\theta, \phi) \boldsymbol{\epsilon}'_1(\mathbf{p}'),\\
    \boldsymbol{\epsilon}_2(\mathbf{p}) &= R(\theta, \phi) \boldsymbol{\epsilon}'_2(\mathbf{p}').
\end{align}
We now apply the rotation matrix to each polarization vector. First, let’s calculate \( \boldsymbol{\epsilon}_1(\mathbf{p}) \):
\begin{align*}
    \boldsymbol{\epsilon}_1(\mathbf{p}) &= 
    \begin{pmatrix}
        \cos \phi & -\sin \phi & 0 \\
        \sin \phi & \cos \phi & 0 \\
        0 & 0 & 1
    \end{pmatrix}
    \begin{pmatrix}
        \cos \theta & 0 & \sin \theta \\
        0 & 1 & 0 \\
        -\sin \theta & 0 & \cos \theta
    \end{pmatrix}
    \begin{pmatrix}
        1 \\
        0 \\
        0
    \end{pmatrix}, \nonumber \\
    &= 
    \begin{pmatrix}
        \cos \phi & -\sin \phi & 0 \\
        \sin \phi & \cos \phi & 0 \\
        0 & 0 & 1
    \end{pmatrix}
    \begin{pmatrix}
        \cos \theta \\
        0 \\
        -\sin \theta
    \end{pmatrix}, \nonumber \\
    &= 
    \begin{pmatrix}
        \cos \theta \cos \phi \\
        \cos \theta \sin \phi \\
        -\sin \theta
    \end{pmatrix}.
\end{align*}
Now, let’s calculate \( \boldsymbol{\epsilon}_2(\mathbf{p}) \):
\begin{align*}
    \boldsymbol{\epsilon}_2(\mathbf{p}) &= 
    \begin{pmatrix}
        \cos \phi & -\sin \phi & 0 \\
        \sin \phi & \cos \phi & 0 \\
        0 & 0 & 1
    \end{pmatrix}
    \begin{pmatrix}
        \cos \theta & 0 & \sin \theta \\
        0 & 1 & 0 \\
        -\sin \theta & 0 & \cos \theta
    \end{pmatrix}
    \begin{pmatrix}
        0 \\
        1 \\
        0
    \end{pmatrix}, \nonumber \\
    &= 
    \begin{pmatrix}
        \cos \phi & -\sin \phi & 0 \\
        \sin \phi & \cos \phi & 0 \\
        0 & 0 & 1
    \end{pmatrix}
    \begin{pmatrix}
        0 \\
        1 \\
        0
    \end{pmatrix}, \nonumber \\
    &= 
    \begin{pmatrix}
        -\sin \phi \\
        \cos \phi \\
        0
    \end{pmatrix}.
\end{align*}

Thus, the two polarization vectors in the unprimed frame are:
\begin{align}
    \boldsymbol{\epsilon}_1(\mathbf{p}) &= 
    \begin{pmatrix}
        \cos \theta \cos \phi \\
        \cos \theta \sin \phi \\
        -\sin \theta
    \end{pmatrix}, \\
    \boldsymbol{\epsilon}_2(\mathbf{p}) &= 
    \begin{pmatrix}
        -\sin \phi \\
        \cos \phi \\
        0
    \end{pmatrix}.
\end{align}
\item [(b)] Let us define the third polarization vector as:
\begin{align}
    \boldsymbol{\epsilon}_3(\mathbf{p}) = \hat{p} = \frac{\mathbf{p}}{|\mathbf{p}|}.
\end{align}
The three polarization vectors $\boldsymbol{\epsilon}_r(\mathbf{p})$ ($r = 1, 2, 3$) form a complete orthonormal basis for the three-dimensional momentum space. Therefore, they obey the completeness relation:
\begin{align}
    \sum_{r=1}^3 \epsilon_i^r(\mathbf{p}) \epsilon_j^r(\mathbf{p}) = \delta_{ij}. \label{eq:completeness-relation}
\end{align}
The vector $\displaystyle \boldsymbol{\epsilon}_3(\mathbf{p}) = \hat{p}$ represents the longitudinal component along the direction of momentum. Therefore, the completeness relation for the transverse polarization vectors is obtained by subtracting the contribution of $\boldsymbol{\epsilon}_3(\mathbf{p})$ from (\ref{eq:completeness-relation}).

We have:
\begin{align*}
    \epsilon_i^3(\mathbf{p}) = \hat{p}_i = \frac{p_i}{|\mathbf{p}|}.
\end{align*}
Now, subtract the longitudinal component from the completeness relation:
\begin{align*}
    \sum_{r=1}^2 \epsilon_i^r(\mathbf{p}) \epsilon_j^r(\mathbf{p}) = \delta_{ij} - \epsilon_i^3(\mathbf{p}) \epsilon_j^3(\mathbf{p}).
\end{align*}
Substitute $\displaystyle\epsilon_i^3(\mathbf{p}) = \frac{p_i}{|\mathbf{p}|}$:
\begin{align*}
    \sum_{r=1}^2 \epsilon_i^r(\mathbf{p}) \epsilon_j^r(\mathbf{p}) = \delta_{ij} - \frac{p_i p_j}{|\mathbf{p}|^2}.
\end{align*}
Thus, the completeness relation for the two transverse polarization vectors is:
\begin{align}
    \sum_{r=1}^2 \epsilon_i^r(\mathbf{p}) \epsilon_j^r(\mathbf{p}) = \delta_{ij} - \frac{p_i p_j}{|\mathbf{p}|^2}.
\end{align}
\end{enumerate}
\bigskip\bigskip\hline\hline\bigskip